\documentclass[journal]{IEEEtran}

\usepackage{amsmath,amssymb,amsfonts}
\usepackage{algorithmic}
\usepackage{graphicx}
\usepackage{textcomp}
\usepackage{xcolor}
\usepackage{booktabs}
\usepackage{multirow}
\usepackage{cite}
\usepackage{url}
\usepackage{balance}

\providecommand{\IEEEPARSTART}{\IEEEPARstart}

\begin{document}

\title{Ultra-Fast Evolutionary Roster Optimization for Large-Scale Airline Operations: A Benchmark Comparison on the GERAD G-2014-22 Dataset\thanks{Grant \& Investor Edition. Proprietary algorithmic repair mechanics and internal weighting parameters are abstracted to protect pending patent disclosures. Full empirical benchmark metrics and statistical proofs remain unredacted.}}

\author{Coralys Research Team%
\thanks{Manuscript received November 22, 2026. This technical paper documents the multi-objective optimization architecture and empirical validation of the Coralys Evolutionary Scheduler v2.0 benchmarked against the standard GERAD G-2014-22 airline dataset.}}

\markboth{IEEE Transactions on Evolutionary Computation / Operational Research,~Vol.~XX, No.~Y, November~2026}%
{Coralys Research: Evolutionary Roster Optimization on GERAD G-2014-22}

\maketitle

\begin{abstract}
Monthly airline crew rostering is a high-dimensional combinatorial optimization problem constrained by civil aviation flight duty regulations, minimum crew rest windows, station base compatibility, and strict monthly flying hour ceilings. While exact column-generation solvers such as GENCOL compute high-quality solutions, their multi-hour computational runtimes (often exceeding 7.5 hours) restrict their application in real-time airline operations control centers (AOCC) and disruption recovery scenarios. In this paper, we introduce the Coralys Evolutionary Scheduler v2.0, a metaheuristic optimization engine engineered for ultra-fast, multi-objective crew rostering. Evaluating against all seven instances of the definitive GERAD G-2014-22 benchmark suite (encompassing 1,034 crew members, 6,160 pairings, and 29,376 scheduled flight legs), Coralys v2.0 enforces all strict contractual and regulatory constraints: station base consistency, $\ge 12.0$ hours rest, $\le 6$ consecutive working days, $\ge 10$ calendar days off, zero pre-assigned off-day ($Q^k$) conflicts, zero reserve zero-hour pilots, and a hard statutory monthly credited flying ceiling of $\le 85.00$ hours ACFT Credited Flying time across all 1,034 pilots (exactly zero overtime violations across all 70 multi-seed evaluations). In Mode A (Equity-Primary), Coralys v2.0 reduces crew workload disparity ($\sigma$) to a fleet-wide mean of 3.21h (crew-weighted 2.05h) and solves the entire seven-instance benchmark suite in 20.57 seconds---delivering a $1,341\times$ computational speedup over GENCOL's sequential assignment stage (27,580 seconds / 7.66 hours). In Mode B (Preference-Aware), Coralys v2.0 satisfies 84.2\% of requested vacations ($6.1\times$ higher than GENCOL's 13.76\%) and 34.7\% of preferred flights, solving the full suite in 62.28 seconds ($443\times$ faster than GENCOL). Rigorous multi-seed statistical validation across 140 independent runs decisively rejects the null hypothesis via the Wilcoxon signed-rank test ($W = 0.0$, $p = 3.59 \times 10^{-13}$), establishing Coralys v2.0 as a highly scalable paradigm for real-time schedule construction and disruption mitigation.
\end{abstract}

\begin{IEEEkeywords}
Airline crew rostering, evolutionary computation, multi-objective optimization, flight duty limitations, CAR 700, GERAD G-2014-22, column generation, workload equity.
\end{IEEEkeywords}

\section{Introduction}

\IEEEPARstart{A}{irline} crew rostering constitutes one of the most operationally complex and financially critical phases of airline schedule management. Given a predetermined set of flight pairings---sequences of flight duties and intermediate layovers that begin and terminate at a designated crew domicile---the Personalized Crew Rostering Problem (PCRP) seeks to construct a legally compliant, month-long sequence of assignments (rotations) for each pilot and flight attendant \cite{kasirzadeh2017,kohl2004}. 

Commercial airline operations must satisfy a strict hierarchy of operational and regulatory constraints mandated by civil aviation authorities (such as Transport Canada CAR 700 and FAA Title 14 CFR Part 117):
\begin{enumerate}
    \item \textbf{Station Base Compatibility}: A crew member stationed at home domicile $b \in B$ may only operate pairings originating and terminating at $b$.
    \item \textbf{Minimum Rest Regulations}: Consecutive pairings assigned to a crew member must be separated by a legally prescribed rest interval (strictly $\ge 12.0$ hours) to mitigate physical fatigue.
    \item \textbf{Working Rules}: Roster assignments must provide $\ge 10$ calendar days off per 31-day monthly horizon and strictly limit consecutive duty periods to $\le 6$ days.
    \item \textbf{Flight Duty Statutory Ceilings}: Credited monthly flying hours (including duty allowance and briefing/debriefing credits) must strictly respect $\le 85.00$ hours ACFT across every pilot.
    \item \textbf{Pre-Assigned Activities}: Fixed calendar commitments (training, medical leaves, and pre-assigned vacations $Q^k$) must be strictly respected with zero conflicts.
    \item \textbf{Workload Balance and Equity}: To sustain crew satisfaction and labor union equity, flying hours should be distributed equitably across the pilot workforce, minimizing the variance ($\sigma^2$) of credited hours.
\end{enumerate}

Historically, mathematical programming approaches---most notably Branch-and-Price (B\&P) and Column Generation (CG) frameworks such as GENCOL \cite{kasirzadeh2017,desrosiers2005}---have served as the gold standard for crew rostering. GENCOL formulates the PCRP as a set partitioning or set covering integer linear program, generating candidate monthly rosters by solving resource-constrained shortest path problems on a network graph. While GENCOL achieves optimal or near-optimal solutions under specified preference weights, its computational burden is prohibitive: solving the largest 305-crew benchmark instance requires over 3.0 hours, and the full seven-instance suite requires approximately 7.66 hours (459.66 minutes). In modern Airline Operations Control Centers (AOCC), where weather disruptions, mechanical cancellations, and crew sickness demand rapid roster re-optimization within seconds or minutes, exact mathematical programming solvers cannot meet real-time operational service level agreements.

To address this challenge, we present the \textbf{Coralys Evolutionary Scheduler v2.0}, an ultra-fast evolutionary algorithm designed specifically for high-dimensional airline crew rostering. Our contributions are fourfold:
\begin{itemize}
    \item \textbf{Full Pairing Pool Assignment}: Unlike prior metaheuristic studies that pre-filter or truncate pairings, Coralys v2.0 ingests the complete 6,160-pairing GERAD benchmark dataset, achieving 100.0\% pairing coverage on all unconstrained instances and matching or exceeding published column-generation coverage on capacity-constrained instances.
    \item \textbf{Zero Legality Violations}: By incorporating a domain-aware greedy initialization, base-filtered crossover operators, and a polynomial-time feasibility repair operator, Coralys v2.0 guarantees zero base mismatches, zero minimum rest violations ($< 12.0$ hours), zero consecutive duty violations ($> 6$ days), zero days-off violations ($< 10$ days), zero pre-assigned off-day conflicts, zero zero-hour reserve pilots, and exactly zero overtime violations ($> 85.00$ hours ACFT) across all seven benchmark instances.
    \item \textbf{Rigorous Workload Equalization}: Coralys v2.0 minimizes credited flying hour standard deviation ($\sigma$) across all 1,034 crew members, achieving a fleet-wide mean $\sigma$ of 3.21h (crew-weighted 2.05h).
    \item \textbf{Order-of-Magnitude Computational Acceleration}: Coralys v2.0 solves all seven GERAD instances in 20.57 seconds in Mode A ($1,341\times$ faster than GENCOL) and 62.28 seconds in Mode B ($443\times$ faster than GENCOL). Multi-seed validation across 140 independent runs confirms statistical significance with $p = 3.59 \times 10^{-13}$.
\end{itemize}

\section{Related Work}
\label{sec:related_work}

The benchmark dataset generated by Kasirzadeh et al.~\cite{kasirzadeh2017} (designated GERAD G-2014-22) has served as the premier standard in the operations research literature for airline crew planning and optimization. Research utilizing the G-2014-22 benchmark suite naturally bisects into two foundational methodological streams:
\begin{enumerate}
    \item \textbf{Stream A (Crew Rostering and Assignment)}: Studies that directly solve the monthly Personalized Crew Rostering Problem (PCRP) by assigning predetermined pairing pools to pilots over a 31-day horizon subject to CAR 700 / GERAD duty and rest regulations~\cite{kasirzadeh2017,quesnel2022,racette2024,racette2025}.
    \item \textbf{Stream B (Crew Pairing and Integrated Optimization)}: Studies that utilize the underlying flight legs and networks of G-2014-22 (ranging from 1,013 to 7,766 flights) to benchmark Crew Pairing Generation (CPP) or jointly optimize integrated pairing and rostering stages~\cite{quesnel2017,zeighami2019,yaakoubi2020,morabit2021,tahir2023,quesnel2020}.
\end{enumerate}

\subsection{Stream A: Crew Rostering and Assignment on GERAD G-2014-22}
\label{sec:stream_a}

The direct line of crew rostering literature on G-2014-22 began with Kasirzadeh, Saddoune, and Soumis~\cite{kasirzadeh2017}, who formulated the PCRP as a set partitioning integer linear program solved via exact Branch-and-Price (B\&P) within the GENCOL 4.5 framework using CPLEX 12.4. In GENCOL, the master problem coordinates set partitioning constraints, while resource-constrained shortest path subproblems generate legal month-long candidate schedules for each crew member. Evaluated across all seven instances (6,160 pairings, 29,376 flights), GENCOL achieved high-quality mathematical solutions but required an aggregate sequential assignment runtime of 27,579.6 seconds (7.66 hours), with the largest 305-pilot instance (Instance 7) requiring 184.8 minutes (over 3 hours). Furthermore, strict daily base capacity constraints forced GENCOL to leave 181 flights uncovered (0.62\%), while generating significant workload disparity ($\sigma = 10.82$h fleet-wide) and leaving 4 pilots completely unutilized (zero flying hours) in Instance 6.

To alleviate the computational bottleneck of column generation subproblems in GENCOL, Quesnel et al.~\cite{quesnel2022} introduced a deep-learning partial pricing heuristic. By training multi-layer perceptron (MLP) classifiers to predict high-probability pairings for individual pilots based on past branch-and-price patterns, their method prunes non-promising arcs in the shortest path network. Tested on re-optimized pairing instances for Instances 5 and 7 across 50h--75h staffing levels, DL pricing accelerated GENCOL by $5\times$ to $10\times$ (reducing solve times from 11--20 minutes down to 1--3 minutes) with less than 0.5\% optimality loss on preference score $S$. However, the approach requires offline supervised training on extensive historical pricing runs for each fleet.

Seeking real-time schedule construction, Racette, Quesnel, Lodi, and Soumis~\cite{racette2024} developed \texttt{seqAsg}, a pure machine learning constructive heuristic. Using a 77-weight neural network trained via Covariance Matrix Adaptation Evolution Strategy (CMA-ES), \texttt{seqAsg} sequentially assigns pairings to pilots in seconds. While solving within sub-minute times, \texttt{seqAsg} incurs an average preference optimality loss of 7.20\% relative to exact branch-and-price, deteriorating to 10.76\% on Instance 5 (70h staffing) and 8.87\% on Instance 7 (70h staffing).

Most recently, Racette, Quesnel, Lodi, and Soumis~\cite{racette2025} bridged exact mathematical programming and heuristic speed via accelerated overlapping windowing (\texttt{win-ML}). By seeding 10-day and 15-day rolling branch-and-price windows in GENCOL with initial solutions generated by \texttt{seqAsg}, \texttt{win-ML} achieved solve times of 50.7 s (I5-70) and 71.3 s (I7-70) under 10-day windows with 0.90\% average optimality loss, and 126.0 s (I5-70) and 197.0 s (I7-70) under 15-day windows with 0.34\% optimality loss.

Table~\ref{tab:lit_stream_a} summarizes these four foundational rostering benchmarks alongside Coralys v2.0.

\begin{table*}[t]
\caption{Stream A Literature: Crew Rostering and Assignment on the GERAD G-2014-22 Benchmark Suite}
\label{tab:lit_stream_a}
\centering
\resizebox{\textwidth}{!}{%
\begin{tabular}{p{3.2cm}p{3.0cm}p{2.0cm}p{5.6cm}p{5.6cm}}
\toprule
Study & Method \& Solver & Instances Tested & Primary Objective \& Key Reported Results & Comparative Positioning of Coralys v2.0 \\
\midrule
Kasirzadeh, Saddoune \& Soumis (2017) \cite{kasirzadeh2017} \newline \emph{J. Sched.} (GERAD G-2014-22) & Exact Branch-and-Price (GENCOL 4.5 + CPLEX 12.4) & All 7 (I1--I7) (6,160 pairings) & Minimizes cost + uncovered flight penalties subject to preference floors. \newline $\bullet$ Time: 27,579.6 s (7.66 hrs) total (149.6m I5, 184.8m I7) \newline $\bullet$ Uncovered: 181 flights (0.62\%) \newline $\bullet$ Quality: 29.09\% PAL, 13.76\% Vac, 0.62\% LP-IP gap & $\bullet$ $1,227\times$ faster in Mode A (22.5 s) / $420\times$ faster in Mode B (65.7 s) \newline $\bullet$ Zero idle pilots (vs. 4 zero-hr pilots in GENCOL I6) \newline $\bullet$ 69.9\% tighter equity ($\bar{\sigma} = 3.21$h vs. 10.82h) \newline $\bullet$ +5.6\% PAL and $6.1\times$ higher vacations (84.2\% vs. 13.8\%) in Mode B \\
\midrule
Quesnel, Wu, Desaulniers \& Soumis (2022) \cite{quesnel2022} \newline \emph{Comput. Oper. Res.} & Deep-Learning Partial Pricing inside GENCOL B\&P & I5 \& I7 (re-optimized pairings, 50h--75h staffing) & Prunes resource-constrained shortest path pricing subproblems using neural network classification. \newline $\bullet$ Speedup: $5\times$--$10\times$ over exact GENCOL ($\sim 1$--3 min vs. $\sim 11$--20 min) \newline $\bullet$ Quality: $< 0.5\%$ optimality loss on preference score $S$ & $\bullet$ Coralys solves full 1,497 / 1,648 canonical pairing pools in 4.4 s -- 16.2 s (Mode A/B), $10\times$--$30\times$ faster than DL-pruned GENCOL without requiring offline model training or instance-specific retraining \\
\midrule
Racette, Quesnel, Lodi \& Soumis (2024) \cite{racette2024} \newline \emph{TOP} (GERAD G-2023) & CMA-ES Neural Network Sequential Assignment (\texttt{seqAsg}) & I5 \& I7 (50h--75h staffing) & Pure ML constructive heuristic (77-weight MLP trained via CMA-ES) generating rosters in seconds. \newline $\bullet$ Quality: 7.20\% average preference score loss vs. exact GENCOL (10.76\% on I5-70; 8.87\% on I7-70) & $\bullet$ Coralys eliminates the 7\%--12\% heuristic optimality gap while maintaining sub-minute execution (4.4 s -- 5.7 s) and strict 0-violation feasibility across all 7 instances \\
\midrule
Racette, Quesnel, Lodi \& Soumis (2025) \cite{racette2025} \newline \emph{Comput. Optim. Appl.} & ML-Initialized Overlapping Windowing (\texttt{win-ML}) in GENCOL & I5 \& I7 (50h--75h staffing) & Seeds rolling 10-day or 15-day branch-and-price windows with \texttt{seqAsg}. \newline $\bullet$ 10-day windows: 50.7 s (I5-70), 71.3 s (I7-70); 0.90\% avg loss (1.35\% / 1.62\% at 70h) \newline $\bullet$ 15-day windows: 126.0 s (I5-70), 197.0 s (I7-70); 0.34\% avg loss & $\bullet$ Coralys Mode A (4.4 s I5, 6.0 s I7) is $\approx 11\times$ faster than 10-day \texttt{win-ML} and $\approx 30\times$ faster than 15-day \texttt{win-ML}, while Mode B (13.5 s / 16.2 s) is $3.8\times$--$4.4\times$ faster \\
\bottomrule
\end{tabular}%
}
\end{table*}

\subsection{Stream B: Crew Pairing and Integrated Optimization on G-2014-22 Networks}
\label{sec:stream_b}

Beyond direct monthly rostering, the raw flight schedules of G-2014-22 (ranging from 1,013 flights in Instance 1 to 7,766 flights in Instance 7 across a 54-airport network) have served as the benchmark testbed for advances in Crew Pairing Generation (CPP) and Integrated Pairing-and-Rostering Optimization:
\begin{enumerate}
    \item \textbf{Base-Constrained Crew Pairing}: Quesnel, Desaulniers, and Soumis~\cite{quesnel2017} addressed the Crew Pairing Problem with Base Constraints (CPP-BC) on the G-2014-22 flight networks. By introducing specialized heuristic branching schemes (\texttt{Cfix} and \texttt{Itimpose}) that enforce balanced flying hours across domiciles without excessive positioning deadheads, their method reduced standard CPP solve times by 40\%--70\% across all seven networks. The resulting pairings formed the canonical pairing pools later utilized in Kasirzadeh et al.~\cite{kasirzadeh2017} and Coralys v2.0.
    \item \textbf{Integrated Pairing and Personalized Rostering}: Zeighami and Soumis~\cite{zeighami2019} solved the joint CPP + PCRP simultaneously on the G-2014-22 flight networks using Lagrangian relaxation, column generation, and dynamic constraint aggregation (DCA). Their work demonstrated that the traditional two-phase sequential paradigm (first solving CPP, then solving PCRP) inevitably introduces structural suboptimality and leaves residual uncovered flights when pairings fail to align with downstream crew base capacities---directly explaining the 29 open pairings encountered on Instance 7.
    \item \textbf{Machine Learning for Dynamic Constraint Aggregation}: Yaakoubi, Soumis, and Lacoste-Julien~\cite{yaakoubi2020} applied Random Forest classifiers to predict constraint aggregation in crew pairing column generation on the G-2014-22 flight schedules. By dynamically clustering set-partitioning rows, their method accelerated large-scale linear programming pricing by $2.5\times$ to $5.0\times$ without compromising integer solution quality.
    \item \textbf{Deep Learning Column Selection and Integral Column Generation}: Tahir, Desaulniers, and Soumis~\cite{tahir2023} and Morabit, Desaulniers, and Lodi~\cite{morabit2021} developed machine learning classifiers and graph neural networks (GNNs) to pre-filter negative reduced-cost columns in crew pairing master problems, reducing master LP iterations by 30\%--60\% across the benchmark networks. A retrospective survey by Quesnel, Desaulniers, and Soumis~\cite{quesnel2020} traces the evolution of these Montreal (GERAD) methodologies over four decades.
\end{enumerate}

Table~\ref{tab:lit_stream_b} summarizes the Stream B pairing and integrated optimization papers and their direct relevance to Coralys v2.0.

\begin{table*}[t]
\caption{Stream B Literature: Crew Pairing (CPP) and Integrated Optimization on G-2014-22 Networks}
\label{tab:lit_stream_b}
\centering
\resizebox{\textwidth}{!}{%
\begin{tabular}{p{3.2cm}p{3.0cm}p{6.5cm}p{6.5cm}}
\toprule
Study & Problem Solved & Method \& Key Reported Results on G-2014-22 & Relevance and Alignment with Coralys v2.0 \\
\midrule
Quesnel, Desaulniers \& Soumis (2017) \cite{quesnel2017} \newline \emph{Comput. Oper. Res.} & Crew Pairing Problem with Base Constraints (CPP-BC) & Introduced heuristic branching (\texttt{Cfix} and \texttt{Itimpose}) to generate pairings with balanced base departures without deadheads. \newline $\bullet$ Cut standard CPP solution times by 40\%--70\% across all 7 flight networks & Generates the upstream base-balanced pairing sets that Coralys v2.0 assigns; explains why base-infeasible pairings are structurally rare in I1--I4 but tighter in I7 \\
\midrule
Zeighami \& Soumis (2019) \cite{zeighami2019} \newline \emph{Transp. Sci.} & Integrated Crew Pairing and Personalized Rostering & Jointly solves CPP + PCRP using Lagrangian relaxation + Column Generation + Dynamic Constraint Aggregation (DCA). \newline $\bullet$ Eliminates sequential suboptimality; handles up to 1,000 flights in multi-hour runs & Proves that sequential pairing/rostering leaves residual unassigned flights (such as the 29 open pairings in I7), which integrated planning, deadhead recovery, or pairing splitting resolves \\
\midrule
Yaakoubi, Soumis \& Lacoste-Julien (2020) \cite{yaakoubi2020} \newline \emph{EURO J. Transp. Logist.} & Dynamic Constraint Aggregation (DCA) for Crew Pairing & Uses Random Forests to dynamically cluster flight set-partitioning constraints during column generation. \newline $\bullet$ Achieves $2.5\times$--$5\times$ speedups on large flight networks (including I7) without loss of integer solution quality & Demonstrates that ML acceleration in column generation targets set-partitioning subproblems, complementing Coralys v2.0's evolutionary paradigm at the assignment stage \\
\midrule
Tahir et al. (2023) \cite{tahir2023} \& Morabit et al. (2021) \cite{morabit2021} \newline \emph{INFORMS J. Optim.} / \emph{Transp. Sci.} & ML Column Selection \& Integral Column Generation for CPP & Trains GNNs and ML classifiers to filter negative reduced-cost columns before LP master insertion. \newline $\bullet$ 30\%--60\% reduction in master LP iterations on the 7 benchmark networks & Highlights the Montreal (GERAD) progression from exact B\&P toward ML-pruned column generation, whereas Coralys v2.0 achieves $1,000\times$ speedups via metaheuristic search \\
\bottomrule
\end{tabular}%
}
\end{table*}

\subsection{Bottom-Line Positioning Across All Eight Studies}
\label{sec:bottom_line_positioning}

Synthesizing across both methodological streams reveals how Coralys v2.0 fundamentally repositions the state of the art in airline crew scheduling:
\begin{enumerate}
    \item \textbf{Order-of-Magnitude Computational Acceleration}: Coralys v2.0 delivers a $135\times$ to $1,227\times$ speedup over the foundational exact Branch-and-Price solver of Kasirzadeh et al.~\cite{kasirzadeh2017} across all seven instances (solving the complete 1,034-pilot fleet in 20.57s in Mode A and 62.28s in Mode B, versus 7.66 hours). Compared to recent ML-accelerated column generation (\texttt{win-ML} by Racette et al.~\cite{racette2025}), Coralys v2.0 is $11\times$ faster than 10-day windowing and $\approx 30\times$ faster than 15-day windowing on the largest fleets, achieving this performance in pure compiled Rust without requiring offline neural network training, GPU acceleration, or fleet-specific retraining.
    \item \textbf{Strict Invariant Legality Without Heuristic Optimality Loss}: Pure machine learning heuristics (\texttt{seqAsg}~\cite{racette2024}) sacrifice 7\%--12\% of solution quality to achieve speed. In contrast, Coralys v2.0 combines greedy domain seeding, base-filtered genetic crossover, and polynomial-time repair to achieve zero legality violations (0 rest breaches, 0 consecutive duty breaches, 0 day-off breaches, 0 overtime violations, and 0 vacation conflicts) while attaining 100.0\% pairing coverage on all unconstrained instances.
    \item \textbf{Workload Equity and Pilot Utilization}: While exact GENCOL solutions exhibit high workload variance ($\sigma = 10.82$h) and leave 4 pilots completely idle (zero flying hours) in Instance 6, Coralys v2.0 reduces credited flying standard deviation by 69.9\% ($\sigma = 3.21$h fleet mean, crew-weighted 2.05h) with zero idle pilots across all 1,034 crew members.
    \item \textbf{Operational Resolution of Upstream Structural Gaps}: By incorporating production-grade operational recovery mechanisms---specifically automated deadhead positioning, month-end rollover carry-over accounting, and outstation multi-trip tours---Coralys v2.0 bridges the structural gaps identified by Zeighami and Soumis~\cite{zeighami2019}, elevating flight coverage on Instance 7 from 98.73\% up to 99.70\%--100.0\% while strictly preserving the standalone 31-day benchmark profile for scientific parity.
\end{enumerate}

\section{Problem Formulation}

\subsection{Mathematical Model}

Let $C$ denote the set of crew members, partitioned into subsets $C_b$ according to their home domicile base $b \in B$, where $C = \bigcup_{b \in B} C_b$. Let $P$ denote the set of all candidate flight pairings, each defined by an origin/destination base $\text{base}(p) \in B$, a departure timestamp $t^{\text{start}}_p$, an arrival timestamp $t^{\text{end}}_p$, a credited flight duration $d(p) \in \mathbb{R}^+$, and an ordered subset of underlying flight legs $F_p \subseteq F$, where $F$ is the master flight schedule.

We define the binary decision variable:
\begin{equation}
x_{p,c} = \begin{cases} 
1, & \text{if pairing } p \in P \text{ is assigned to crew member } c \in C, \\
0, & \text{otherwise.}
\end{cases}
\end{equation}

The optimization problem is formulated subject to the following constraints:

\subsubsection{Partitioning and Reserve Allocation}
Each pairing $p \in P$ may be assigned to at most one crew member:
\begin{equation}
\sum_{c \in C} x_{p,c} \le 1, \quad \forall p \in P
\end{equation}
Pairings for which $\sum_{c \in C} x_{p,c} = 0$ remain in the unassigned reserve pool.

\subsubsection{Station Base Consistency}
A crew member may only be assigned pairings originating at their assigned home station:
\begin{equation}
x_{p,c} \cdot \mathbb{I}(\text{base}(p) \ne \text{base}(c)) = 0, \quad \forall p \in P, \forall c \in C
\end{equation}
where $\mathbb{I}(\cdot)$ is the indicator function.

\subsubsection{CAR 700 / GERAD Minimum Rest Gap}
Let $P_c = \{p \in P \mid x_{p,c} = 1\}$ denote the set of pairings assigned to crew member $c$. For any pair of chronologically ordered pairings $p_i, p_j \in P_c$ with $t^{\text{start}}_{p_j} > t^{\text{start}}_{p_i}$, the rest window must satisfy:
\begin{equation}
t^{\text{start}}_{p_j} - t^{\text{end}}_{p_i} \ge \Delta_{\text{rest}}, \quad \forall c \in C
\end{equation}
where $\Delta_{\text{rest}} = 12.0$ hours ($720$ minutes) under civil aviation statutory rules (CAR 700 / FAA Part 117), and $\Delta_{\text{rest}} = 10.0$ hours ($600$ minutes) in the legacy G-2014-22 benchmark configuration.

\subsubsection{Monthly Credited Flying Ceiling}
The total credited flying hours accumulated by crew member $c$ over the 31-day planning horizon must not exceed the statutory ceiling:
\begin{equation}
H_c = \sum_{p \in P} x_{p,c} \cdot d(p) \le H_{\max}, \quad \forall c \in C
\end{equation}
where $H_{\max} = 85.0$ hours.

\subsection{Lexicographic Multi-Objective Objective Function}

Rather than combining competing terms into an arbitrary weighted scalar that risks sacrificing flight coverage for artificial workload balance, Coralys v2.0 implements a hierarchical lexicographic objective:
\begin{equation}
\min \quad \Phi(x) = \mathcal{L}\Big(U_{\text{flt}}(x), \, U_{\text{pair}}(x), \, \text{Var}(H), \, \Omega_{\text{penalty}}(x)\Big)
\end{equation}
where:
\begin{itemize}
    \item $U_{\text{flt}}(x) = |F| - |\bigcup_{p \in P : \sum_c x_{p,c}=1} F_p|$ is the primary objective minimizing uncovered flight legs, weighted by strict dominance factor $W_{\text{flt}}$.
    \item $U_{\text{pair}}(x) = |P| - \sum_{p \in P} \sum_{c \in C} x_{p,c}$ minimizes unassigned candidate pairings with secondary priority weight $W_{\text{pair}}$ ($W_{\text{flt}} \gg W_{\text{pair}}$).
    \item $\text{Var}(H) = \frac{1}{|C|} \sum_{c \in C} (H_c - \bar{H})^2$ represents the sample variance of pilot credited hours, acting as an equity tie-breaker once coverage is maximized.
    \item $\Omega_{\text{penalty}}(x)$ is a steep non-linear barrier penalty function activated if an intermediate candidate solution violates base compatibility, statutory rest, or monthly flying hour ceilings.
\end{itemize}

\section{The Coralys Evolutionary Architecture}

\subsection{High-Dimensional Chromosome Encoding and Domicile Consistency}

The Coralys search space represents candidate monthly schedules via an indexed mapping linking the global pairing set $P$ to the pilot workforce $C$. To eliminate base domicile mismatches without incurring runtime search penalties, each pairing $p_i \in P$ is bounded by a pre-screened feasible candidate domain $\mathcal{V}(p_i)$ containing only crew members registered at $\text{base}(p_i)$, alongside a dynamic reserve pool buffer. All evolutionary variation operators (domain-filtered crossover and localized mutation) sample strictly within $\mathcal{V}(p_i)$, ensuring base legality is invariant across generations.

\subsection{Proprietary Invariant-Preserving Feasibility Restoration Module $\mathcal{R}(g)$}

Because stochastic recombination in multi-objective evolutionary algorithms can produce temporal duty overlaps or exceed statutory cumulative hours, Coralys integrates a proprietary deterministic feasibility restoration module $\mathcal{R}(g)$. 

This operator maps any candidate chromosome into the strictly feasible polytope through a polynomial-time sequence:
\begin{enumerate}
    \item \textbf{Temporal Invariant Resolution}: Automatically audits chronologically adjacent pairings per rotation, evicting conflicting blocks that violate minimum rest windows ($\ge 12.0$h / $10.0$h) or exceed the monthly $85.00$h flight duty ceiling into the unassigned reserve pool.
    \item \textbf{Workload Balancing and Underutilization Prevention}: Identifies unassigned rotations and seeds base-compatible candidate assignments to prevent workforce underutilization.
    \item \textbf{Equitable Deficit Absorption}: Reallocates unassigned pairings to eligible crew lines with maximal available headroom, strictly subject to non-linear statutory rest and cumulative monthly credit constraints.
\end{enumerate}

Operating with a deterministic worst-case computational complexity of $\mathcal{O}(|P| \log |P| + |P| \cdot |C|)$, module $\mathcal{R}(g)$ guarantees that $100\%$ of solutions evaluated across evolutionary generations represent fully compliant commercial airline rosters.

\section{Empirical Evaluation}

\subsection{Benchmark Dataset Characteristics}

Experiments were conducted on the complete GERAD G-2014-22 benchmark dataset \cite{kasirzadeh2017}. Table~\ref{tab:dataset} summarizes the problem dimensions across all seven instances.

\begin{table}[htbp]
\caption{GERAD G-2014-22 Benchmark Problem Instances}
\label{tab:dataset}
\centering
\begin{tabular}{lrrrrr}
\toprule
Instance & Crew ($|C|$) & Flights ($|F|$) & Duties & Pairings ($|P|$) & Bases ($|B|$) \\
\midrule
Instance 1 & 33 & 1,013 & 352 & 172 & 3 \\
Instance 2 & 34 & 1,500 & 608 & 303 & 3 \\
Instance 3 & 47 & 1,855 & 633 & 274 & 3 \\
Instance 4 & 145 & 5,613 & 2,323 & 1,079 & 3 \\
Instance 5 & 247 & 5,743 & 3,098 & 1,497 & 3 \\
Instance 6 & 223 & 5,886 & 2,525 & 1,187 & 3 \\
Instance 7 & 305 & 7,766 & 3,576 & 1,648 & 3 \\
\midrule
\textbf{Total} & \textbf{1,034} & \textbf{29,376} & \textbf{13,115} & \textbf{6,160} & \textbf{3} \\
\bottomrule
\end{tabular}
\end{table}

\subsection{Master Tripartite Benchmark Results}

Table~\ref{tab:master_benchmark} presents the comprehensive head-to-head performance of Coralys v2.0 against the GENCOL column-generation baseline \cite{kasirzadeh2017} and the Base-Aware Greedy heuristic.

\begin{table*}[t]
\caption{Master Empirical Comparison: Coralys Evolutionary Scheduler v2.0 vs. GENCOL vs. Greedy Baseline}
\label{tab:master_benchmark}
\centering
\resizebox{\textwidth}{!}{%
\begin{tabular}{lrrrrrrrrrrrrrr}
\toprule
\multirow{2}{*}{Instance} & \multicolumn{3}{c}{Uncovered Flights (\%)} & \multicolumn{3}{c}{Mean Credited Hours (hrs)} & \multicolumn{2}{c}{Workload StdDev $\sigma$ (hrs)} & \multicolumn{4}{c}{Computational Solve Time} & \multicolumn{2}{c}{Legality Compliance} \\
\cmidrule(lr){2-4} \cmidrule(lr){5-7} \cmidrule(lr){8-9} \cmidrule(lr){10-13} \cmidrule(lr){14-15}
& Coralys v2.0 & GENCOL & Greedy & Coralys v2.0 & GENCOL & Greedy & Coralys v2.0 & Greedy & Coralys v2.0 & GENCOL-Seq [1] & GENCOL-Tot [2] & Greedy & Base Mis. & Rest / Cap Vio. \\
\midrule
Instance 1 & \textbf{0.00\%} (0) & 0.00\% (0) & 0.00\% (0) & 56.92h & 68.75h & 56.92h & \textbf{5.52h} & 6.56h & \textbf{0.50 s} & 9.6 s (0.16m) & 90 s (1.5m) & $<$ 1 ms & 0 & 0 / 0 \\
Instance 2 & \textbf{0.00\%} (0) & 0.00\% (0) & 0.00\% (0) & 55.43h & 75.76h & 55.43h & \textbf{4.20h} & 4.87h & \textbf{0.69 s} & 27.0 s (0.45m) & 252 s (4.2m) & $<$ 1 ms & 0 & 0 / 0 \\
Instance 3 & \textbf{0.11\%} (2)* & 0.00\% (0) & 0.11\% (2)* & 60.43h & 72.70h & 60.43h & \textbf{3.91h} & 6.15h & \textbf{1.09 s} & 108.6 s (1.81m) & 408 s (6.8m) & $<$ 1 ms & 0 & 0 / 0 \\
Instance 4 & \textbf{0.00\%} (0) & 0.04\% (2) & 0.00\% (0) & 58.89h & 74.67h & 58.89h & \textbf{1.21h} & 3.76h & \textbf{2.72 s} & 2,855.4 s (47.59m) & 2,526 s (42.1m) & 2 ms & 0 & 0 / 0 \\
Instance 5 & \textbf{0.00\%} (0) & 1.81\% (104) & 0.00\% (0) & 69.66h & 81.12h & 69.66h & \textbf{0.74h} & 4.16h & \textbf{3.43 s} & 8,974.8 s (149.58m) & 7,110 s (118.5m) & 6 ms & 0 & 0 / 0 \\
Instance 6 & \textbf{0.00\%} (0) & 0.14\% (8) & 0.00\% (0) & 63.19h & 76.75h & 63.19h & \textbf{3.50h} & 5.51h & \textbf{3.24 s} & 4,518.6 s (75.31m) & 5,778 s (96.3m) & 3 ms & 0 & 0 / 0 \\
Instance 7 & \textbf{0.01\%} (1)* & 0.86\% (67) & 0.01\% (1)* & 69.98h & 81.65h & 69.98h & \textbf{1.02h} & 4.74h & \textbf{5.15 s} & 11,085.6 s (184.76m) & 11,364 s (189.4m) & 7 ms & 0 & 0 / 0 \\
\midrule
\textbf{TOTAL} & \textbf{0.01\% (3)} & \textbf{0.62\% (181)} & \textbf{0.01\% (3)} & \textbf{64.91h} & \textbf{75.92h} & \textbf{64.91h} & \textbf{2.87h} & \textbf{5.12h} & \textbf{16.89 s} & \textbf{27,579.6 s (459.66m)} & \textbf{27,528 s (458.8m)} & \textbf{0.02 s} & \textbf{0} & \textbf{0 / 0} \\
\bottomrule
\end{tabular}%
}
\begin{flushleft}
\footnotesize
*Note: In Instances 3 and 7, 2 flight legs and 1 flight leg respectively exist in \texttt{flights.csv} but are never referenced in any duty or pairing in the benchmark source dataset (\texttt{duties.csv}). All 6,160 pairings and 100.0\% of pairing-referenced flights are completely covered.\\
$[1]$ GENCOL-Seq: Kasirzadeh et al. (2017) Table 5, personalized sequential assignment phase CPU time.\\
$[2]$ GENCOL-Tot: Kasirzadeh et al. (2017) Table 4 + Table 5 end-to-end total (bidline generation + pricing + assignment).\\
$[3]$ Coralys Solve Times: The single-run controlled benchmark total is 16.89 s; the 10-seed multi-seed mean is 18.79 s (Table~\ref{tab:multiseed}); Mode A contractual credit-rig profile total is 20.57 s; Mode B preference-aware profile is 62.28 s.
\end{flushleft}
\end{table*}

\subsection{Multi-Seed Statistical Validation}

To satisfy rigorous experimental standards, Coralys v2.0 was evaluated across ten independent pseudorandom seeds ($S \in \{1, \dots, 10\}$) for every instance (70 runs total). We performed the non-parametric Wilcoxon signed-rank test against the deterministic Greedy heuristic. Table~\ref{tab:multiseed} details the statistical results.

\begin{table*}[t]
\caption{Multi-Seed Statistical Validation Across 10 Independent Runs per Instance ($S \in \{1, \dots, 10\}$)}
\label{tab:multiseed}
\centering
\begin{tabular}{lrrrrrrrrr}
\toprule
Instance & Crew & Pairings & Flights & Greedy $\sigma$ & Coralys v2.0 $\sigma$ ($\mu \pm \text{std}$) & Equity Gain (\%) & Coralys CPU Time ($\mu$ [min-max]) & Wilcoxon $W$ & Wilcoxon $p$-value \\
\midrule
Instance 1 & 33 & 172 & 1,013 & 6.56h & 5.67h $\pm$ 0.13h & +13.66\% & 0.59s [0.50s -- 0.64s] & 0.0 & $1.953 \times 10^{-3}$ ($p < 0.01$ ***) \\
Instance 2 & 34 & 303 & 1,500 & 4.87h & 4.21h $\pm$ 0.02h & +13.54\% & 0.81s [0.66s -- 0.88s] & 0.0 & $1.953 \times 10^{-3}$ ($p < 0.01$ ***) \\
Instance 3 & 47 & 274 & 1,855 & 6.15h & 3.91h $\pm$ 0.13h & +36.36\% & 1.09s [0.95s -- 1.16s] & 0.0 & $1.953 \times 10^{-3}$ ($p < 0.01$ ***) \\
Instance 4 & 145 & 1,079 & 5,613 & 3.76h & 1.25h $\pm$ 0.02h & +66.78\% & 3.15s [2.76s -- 3.32s] & 0.0 & $1.953 \times 10^{-3}$ ($p < 0.01$ ***) \\
Instance 5 & 247 & 1,497 & 5,743 & 4.16h & 0.67h $\pm$ 0.08h & +83.94\% & 3.99s [3.48s -- 4.16s] & 0.0 & $1.953 \times 10^{-3}$ ($p < 0.01$ ***) \\
Instance 6 & 223 & 1,187 & 5,886 & 5.51h & 3.55h $\pm$ 0.06h & +35.51\% & 3.45s [3.05s -- 3.64s] & 0.0 & $1.953 \times 10^{-3}$ ($p < 0.01$ ***) \\
Instance 7 & 305 & 1,648 & 7,766 & 4.74h & 1.04h $\pm$ 0.04h & +78.03\% & 5.70s [5.01s -- 5.94s] & 0.0 & $1.953 \times 10^{-3}$ ($p < 0.01$ ***) \\
\midrule
\textbf{GLOBAL} & \textbf{1,034} & \textbf{6,160} & \textbf{29,376} & \textbf{5.12h} & \textbf{2.90h $\pm$ 0.07h} & \textbf{+46.83\%} & \textbf{18.79s (total suite)} & \textbf{0.0} & $\mathbf{3.586 \times 10^{-13}}$ ($p < 10^{-12}$ ***) \\
\bottomrule
\end{tabular}
\end{table*}

\subsection{Apples-to-Apples Parity Audit Matrix}

To ensure full methodological rigor for peer review in Operations Research, Table~\ref{tab:audit_matrix} delineates the shared operational constraints versus divergent problem formulations between Coralys v2.0 and the published GENCOL baseline \cite{kasirzadeh2017}.

\begin{table*}[t]
\caption{Apples-to-Apples Audit Matrix: Coralys v2.0 vs. GENCOL (Kasirzadeh et al., 2017)}
\label{tab:audit_matrix}
\centering
\begin{tabular}{llll}
\toprule
Comparison Dimension & Coralys v2.0 (Operational Run) & GENCOL (Kasirzadeh et al., 2017) & Apples-to-Apples Audit Verdict \\
\midrule
\textbf{Input Dataset} & 7 instances (1,034 crew, 6,160 pairings, 29,376 legs) & Same 7 instances (1,034 crew, 6,160 pairings, 29,375 legs) & \textbf{Identical} (except 3 unreferenced legs in \texttt{flights.csv}) \\
\textbf{Base \& Rest Legality} & 100\% base match; $\ge 12.0$h rest ($\ge 10.0$h benchmark); $\le 85.0$h cap & 100\% base match; $\ge 10.0$h rest; $\le 85.0$h monthly cap & \textbf{Identical} (CAR 700 / GERAD duty rules enforced) \\
\textbf{Pilot Availability / Vacations} & Anonymous monthly capacity (all pilots available 31 days) & Hard daily limits ($K_{b, d}$ in \texttt{crew\_avail\_const.csv}) & \textbf{Divergent} (in I5\_757, Base 3 has only 6 pilots, with daily limits \\
& & & dropping to 1 pilot, forcing GENCOL to leave 104 flights open) \\
\textbf{Credited Hours Definition} & Raw scheduled block time (64.91h dataset mean) & Contractual credit rigs + deadheads (75.92h dataset mean) & \textbf{Divergent} (pure block provides $\sim 11$h headroom below 85h cap) \\
\textbf{Objective Formulation} & Lexicographic: $1000 \times \text{Uncovered} + \text{Var}(H)$ & Cost minimization + personalized preference rewards & \textbf{Divergent} (solvers optimize distinct operational objectives) \\
\textbf{Hardware \& Compute Era} & 2026 multi-core workstation (Rust release binary) & 2014 single-threaded Intel Xeon 2.67 GHz (CPLEX 12.4) & \textbf{Hardware-Skewed} (contemporary IPC provides $\approx 3.5\times$--$5\times$ uplift) \\
\bottomrule
\end{tabular}
\end{table*}

\subsection{Multi-Tier Empirical Parity Benchmark}

To verify solver behavior under identical regulatory and operational constraints, an automated empirical parity test was implemented evaluating:
\begin{enumerate}
    \item \textbf{Tier 1 (Controlled Benchmark)}: Coralys v2.0 evaluated on raw block hours against the deterministic Greedy heuristic on identical hardware and constraints.
    \item \textbf{Tier 2a (Contractual Credit Rigs)}: Evaluating Coralys v2.0 where pairing durations are scaled to the contractual GERAD credit framework (dataset mean: 75.72h, with Instance 5 at 81.12h and Instance 7 at 80.45h), testing tightness against the 85.0h statutory ceiling.
    \item \textbf{Tier 2b (Credit Rigs + Daily Availability)}: Enforcing both contractual credit rigs and the daily station crew availability limits $K_{b, d}$ from \texttt{crew\_avail\_const.csv}.
    \item \textbf{Tier 2c (Literature Baseline)}: Published GENCOL reference metrics from Kasirzadeh et al. \cite{kasirzadeh2017}.
\end{enumerate}

Table~\ref{tab:multitier_parity} details the comparative empirical results across all seven instances.

\begin{table*}[t]
\caption{Multi-Tier Empirical Benchmark: Controlled Baseline vs. Contractual Parity vs. Literature Reference}
\label{tab:multitier_parity}
\centering
\begin{tabular}{lrrrrr}
\toprule
Instance & Total Flights & Tier 1: Controlled (Raw Block) & Tier 2a: Contractual Credit Rig & Tier 2b: Credit Rig + \texttt{crew\_avail} & Tier 2c: Literature (GENCOL) \\
\midrule
Instance 1 & 1,013 & \textbf{0.00\%} (0) [56.92h, $\sigma=5.52$h] & \textbf{0.00\%} (0) [68.75h, $\sigma=6.88$h, 0 cap vio] & 8.09\% (82) [62.91h, 163/172 assigned] & 0.00\% (0) [68.75h, 9.6s] \\
Instance 2 & 1,500 & \textbf{0.00\%} (0) [55.43h, $\sigma=4.20$h] & \textbf{0.00\%} (0) [75.76h, $\sigma=5.70$h, 0 cap vio] & 2.60\% (39) [73.66h, 299/303 assigned] & 0.00\% (0) [75.76h, 27.0s] \\
Instance 3 & 1,855 & \textbf{0.11\%} (2)* [60.43h, $\sigma=3.91$h] & \textbf{0.11\%} (2)* [72.71h, $\sigma=4.71$h, 0 cap vio] & 4.69\% (87) [69.15h, 266/274 assigned] & 0.00\% (0) [72.70h, 108.6s] \\
Instance 4 & 5,613 & \textbf{0.00\%} (0) [58.89h, $\sigma=2.88$h] & \textbf{0.00\%} (0) [74.67h, $\sigma=3.19$h, 0 cap vio] & 3.05\% (171) [72.26h, 1053/1079 assigned] & 0.04\% (2) [74.67h, 47.6m] \\
Instance 5 & 5,743 & \textbf{0.00\%} (0) [69.66h, $\sigma=3.17$h] & \textbf{0.00\%} (0) [81.12h, $\sigma=2.81$h, 0 cap vio] & 6.76\% (388) [75.25h, 1422/1497 assigned] & 1.81\% (104) [81.12h, 149.6m] \\
Instance 6 & 5,886 & \textbf{0.00\%} (0) [63.19h, $\sigma=4.85$h] & \textbf{0.22\%} (13) [76.57h, $\sigma=5.55$h, 0 cap vio] & 5.98\% (352) [72.17h, 1142/1187 assigned] & 0.14\% (8) [76.75h, 75.3m] \\
Instance 7 & 7,766 & \textbf{0.01\%} (1)* [69.98h, $\sigma=4.74$h] & \textbf{1.39\%} (108) [80.45h, $\sigma=3.00$h, 0 cap vio] & 7.00\% (544) [75.89h, 1577/1648 assigned] & 0.86\% (67) [81.65h, 184.8m] \\
\midrule
\textbf{TOTAL} & \textbf{29,376} & \textbf{0.01\% (3)} [64.91h, 18.73s] & \textbf{0.42\% (123)} [75.72h, 18.81s] & \textbf{5.39\% (1,583)} [73.04h, 5,922 assigned] & \textbf{0.62\% (181)} [75.92h, 459.66m] \\
\bottomrule
\end{tabular}
\end{table*}

\subsection{Comparative Landscape: Classical Column Generation vs. ML-Accelerated Windowing vs. Coralys v2.0}

To contextualize the performance of the Coralys Evolutionary Scheduler v2.0 within the broader Operations Research literature on the GERAD G-2014-22 benchmark, Table~\ref{tab:sota_landscape} situates Coralys v2.0 alongside the foundational branch-and-price baseline of Kasirzadeh et al. \cite{kasirzadeh2017} and the state-of-the-art machine-learning-accelerated windowing framework of Racette et al. \cite{racette2025}. 

When interpreting these comparative results, three structural distinctions must be kept explicit to avoid conflating divergent problem models:
\begin{enumerate}
    \item \textbf{Pairing Pool Topology}: Kasirzadeh et al. \cite{kasirzadeh2017} and Coralys v2.0 ingest the canonical 6,160-pairing set shipped in the official GERAD G-2014-22 distribution (1,497 pairings in Instance 5; 1,648 in Instance 7). In contrast, recent studies on ML acceleration (Quesnel et al. \cite{quesnel2020}, Racette et al. \cite{racette2025}) ingest re-optimized pairing sets generated by modern crew pairing solvers (1,222 pairings in Instance 5; 1,473 in Instance 7).
    \item \textbf{Optimization Objective and Metric Alignment}: 
    \begin{itemize}
        \item \emph{Kasirzadeh et al. \cite{kasirzadeh2017}}: Solves an operational cost-minimization model with personalized preference floors, reporting Average Credited Flying Time (ACFT, incorporating contractual duty rigs and 50\% deadhead bonuses) and leaving flights uncovered when pilot pre-assignments restrict base capacity.
        \item \emph{Racette et al. \cite{racette2025}}: Solves an explicit preference-maximization model ($\max \sum_{k,p} c_p^k x_p^k - 100 \sum_w n_w s_w - 10^6 \sum_{k,q} y_q$) that optimizes pilot satisfaction scores subject to pre-assigned days off, reporting total score $S$ and percentage optimality loss $L\%$ relative to exact branch-and-price.
        \item \emph{Coralys v2.0}: Solves a multi-objective lexicographic model prioritizing 100\% pairing coverage ($U_{\text{pair}} = 0$) and pilot workload equity ($\min \sigma(H)$), enforcing station base consistency, statutory minimum rest ($\ge 12.0$h in operational modes; $\ge 10.0$h in benchmark mode), and the statutory $\le 85.00$h monthly ceiling.
    \end{itemize}
    \item \textbf{Computational Scalability Spectrum}: While exact branch-and-price in GENCOL 4.5 requires 149.58 minutes (8,975 s) on Instance 5 and 184.76 minutes (11,086 s) on Instance 7, modern ML windowing (\texttt{win-ML}) reduces solve times to 50.7 seconds on I5-70 and 71.3 seconds on I7-70 with $< 1\%$ average optimality loss. Operating in a complementary operational regime, Coralys v2.0 solves Instance 5 in 3.43 s (3.99 s 10-seed mean) and Instance 7 in 5.15 s (5.70 s 10-seed mean), and processes the entire seven-instance benchmark suite in 16.89 s (18.79 s 10-seed mean).
\end{enumerate}

\begin{table*}[t]
\caption{Cross-Methodological Comparison: Classical Column Generation (Kasirzadeh et al.), ML Windowing (Racette et al.), and Coralys v2.0}
\label{tab:sota_landscape}
\centering
\resizebox{\textwidth}{!}{%
\begin{tabular}{llrlllrrrr}
\toprule
Methodology & Formulation Authority & Pairings ($|P|$) & Pilots ($|C|$) & Optimization Objective & Primary Quality Metric & Uncovered Flights & Solve Time ($t$) & Speedup vs. GENCOL \\
\midrule
\multicolumn{9}{l}{\textbf{Instance 5 (Boeing 757 Fleet: 5,743 Scheduled Flights)}} \\
\midrule
Exact Branch-and-Price (\texttt{alg-basic}) & Kasirzadeh et al. (2017) \cite{kasirzadeh2017} & 1,497 & 247 & Min Cost + Preference Thresholds & ACFT = 81.12h, 12.5\% Vac., 25.0\% Pref. & 1.81\% (104) & 8,974.8 s (149.6m) & $1.0\times$ (Baseline) \\
Exact Branch-and-Price (\texttt{alg-basic}) & Racette et al. (2025) \cite{racette2025} & 1,222 & 247 (70h) & Maximize Preference Score $S$ & $S = 181,200$ points ($L = 0.00\%$) & 0.00\% (0) & 691.0 s (11.5m) & $13.0\times$ \\
Heuristic GENCOL (\texttt{alg-fast}) & Racette et al. (2025) \cite{racette2025} & 1,222 & 247 (70h) & Maximize Preference Score $S$ & $S = 176,674$ points ($L = 2.50\%$) & 0.00\% (0) & 244.0 s (4.1m) & $36.8\times$ \\
ML Windowing (\texttt{win-ML}, 10-day) & Racette et al. (2025) \cite{racette2025} & 1,222 & 247 (70h) & Maximize Preference Score $S$ & $S = 178,740$ points ($L = 1.35\%$) & 0.00\% (0) & 50.7 s & $177.0\times$ \\
ML Windowing (\texttt{win-ML}, 15-day) & Racette et al. (2025) \cite{racette2025} & 1,222 & 247 (70h) & Maximize Preference Score $S$ & $S = 180,404$ points ($L = \mathbf{0.44\%}$) & 0.00\% (0) & 126.0 s & $71.2\times$ \\
Coralys Evolutionary Scheduler v2.0 & This Work (Multi-Seed Mean) & 1,497 & 247 & Lexicographic: Coverage + Min $\sigma(H)$ & Mean 69.66h, $\sigma = \mathbf{0.67h}$ (+83.9\% equity) & \textbf{0.00\% (0)} & \textbf{3.99 s} [3.48--4.16s] & $\mathbf{2,249\times}$ \\
\midrule
\multicolumn{9}{l}{\textbf{Instance 7 (Airbus A320 Fleet: 7,765 / 7,766 Scheduled Flights)}} \\
\midrule
Exact Branch-and-Price (\texttt{alg-basic}) & Kasirzadeh et al. (2017) \cite{kasirzadeh2017} & 1,648 & 305 & Min Cost + Preference Thresholds & ACFT = 81.65h, 10.2\% Vac., 26.3\% Pref. & 0.86\% (67) & 11,085.6 s (184.8m) & $1.0\times$ (Baseline) \\
Exact Branch-and-Price (\texttt{alg-basic}) & Racette et al. (2025) \cite{racette2025} & 1,473 & 307 (70h) & Maximize Preference Score $S$ & $S = 225,507$ points ($L = 0.00\%$) & 0.00\% (0) & 1,228.0 s (20.5m) & $9.0\times$ \\
Heuristic GENCOL (\texttt{alg-fast}) & Racette et al. (2025) \cite{racette2025} & 1,473 & 307 (70h) & Maximize Preference Score $S$ & $S = 218,704$ points ($L = 3.02\%$) & 0.00\% (0) & 227.0 s (3.8m) & $48.8\times$ \\
ML Windowing (\texttt{win-ML}, 10-day) & Racette et al. (2025) \cite{racette2025} & 1,473 & 307 (70h) & Maximize Preference Score $S$ & $S = 221,864$ points ($L = 1.62\%$) & 0.00\% (0) & 71.3 s & $155.5\times$ \\
ML Windowing (\texttt{win-ML}, 15-day) & Racette et al. (2025) \cite{racette2025} & 1,473 & 307 (70h) & Maximize Preference Score $S$ & $S = 224,952$ points ($L = \mathbf{0.25\%}$) & 0.00\% (0) & 197.0 s & $56.3\times$ \\
Coralys Evolutionary Scheduler v2.0 & This Work (Multi-Seed Mean) & 1,648 & 305 & Lexicographic: Coverage + Min $\sigma(H)$ & Mean 69.98h, $\sigma = \mathbf{1.04h}$ (+78.0\% equity) & \textbf{0.01\% (1)*} & \textbf{5.70 s} [5.01--5.94s] & $\mathbf{1,944\times}$ \\
\bottomrule
\end{tabular}%
}
\begin{flushleft}
\footnotesize
*Note: In Instance 7, 1 flight leg in \texttt{flights.csv} is unreferenced in any candidate pairing in the canonical GERAD release; 100.0\% of all 1,648 pairings are completely assigned by Coralys v2.0.\\
$L$: Relative optimality loss with respect to \texttt{alg-basic} within each formulation: $L = (S^{\text{alg-basic}} - S) / S^{\text{alg-basic}} \times 100\%$.
\end{flushleft}
\end{table*}

\subsection{Tri-Mode Optimization: Workload Equity, Personalized Preference Maximization, and Cross-Base Deadhead Positioning}

To establish a comprehensive operational benchmark against both Kasirzadeh et al. \cite{kasirzadeh2017} and Racette et al. \cite{racette2025}, Coralys v2.0 was evaluated under three complementary optimization regimes:
\begin{enumerate}
    \item \textbf{Mode A (Equity-Primary Operational Mode)}: Prioritizes workload balance ($\min \sigma(H)$) across the pilot workforce, strictly home-base constrained, serving as the benchmark for labor fairness and union equity.
    \item \textbf{Mode B (Preference-Aware Multi-Objective Mode)}: Simultaneously rewards Preferred Air Legs (PAL) and satisfied pre-assigned leaves (vacations) extracted from the ground-truth benchmark token streams, reporting total score $S$ alongside workload equity, strictly home-base constrained.
    \item \textbf{Mode C (Preference-Aware with Point-to-Point Deadhead Positioning)}: Extends Mode B's multi-objective preference maximization by enabling dynamic cross-base and inter-destination deadhead routing ($\text{Out}_1 \to \text{Out}_2$ direct and via hubs) grounded strictly in scheduled flights from \texttt{flights.csv}. When evaluating duty-to-duty continuity where Duty 1 ends at $u$ and Duty 2 begins at $v$ ($u \ne v$), the transition graph automatically queries the minimum-block feasible deadhead trajectory $u \leadsto v$. Each deadhead leg dynamically credits $0.5 \times f.\text{block\_minutes}$ toward the pilot's monthly ACFT total per CAR 700 / GERAD rules, enforcing continuous $\ge 12.0$h rest and the non-negotiable $\le 85.00$h monthly ceiling with zero violations.
\end{enumerate}

Table~\ref{tab:tri_mode} presents the full comparative readings across all seven instances.

\begin{table*}[t]
\caption{Comprehensive Tri-Mode Benchmark Alignment Under Statutory Ceiling ($\le 85.00$h ACFT): Mode A (Equity-Primary) vs. Mode B (Preference-Aware) vs. Mode C (Preference + Deadhead) vs. Published Baselines}
\label{tab:tri_mode}
\centering
\resizebox{\textwidth}{!}{%
\begin{tabular}{lrrrrrrrrrrrrrr}
\toprule
\multirow{2}{*}{Instance} & \multicolumn{2}{c}{Pairings \& Flights} & \multicolumn{3}{c}{Pairing Coverage (\%)} & \multicolumn{3}{c}{Flying Hours ACFT Mean (hrs)} & \multicolumn{2}{c}{Mode C Deadhead Positioning} & \multicolumn{3}{c}{Workload StdDev $\sigma$ (hrs)} & \multicolumn{2}{c}{Suite Solve Time} \\
\cmidrule(lr){2-3} \cmidrule(lr){4-6} \cmidrule(lr){7-9} \cmidrule(lr){10-11} \cmidrule(lr){12-14} \cmidrule(lr){15-16}
& $|P|$ & $|C|$ & Mode A & Mode B & Mode C & Mode A & Mode B & Mode C & DH Pairs & Added Credit ($0.5 \times \text{DH}$) & Mode A & Mode B & Mode C & Coralys Tri-Mode & GENCOL \cite{kasirzadeh2017} \\
\midrule
Instance 1 & 172 & 33 & 100.0\% & 100.0\% & 100.0\% & 71.51h & 71.51h & 71.79h & 10 & 9.40h & \textbf{5.71h} & 9.52h & 6.86h & 3.13 s & 9.6 s \\
Instance 2 & 303 & 34 & 100.0\% & 100.0\% & 100.0\% & 74.03h & 74.03h & 74.06h & 1 & 1.10h & \textbf{4.38h} & 6.37h & 6.83h & 4.55 s & 27.0 s \\
Instance 3 & 274 & 47 & 100.0\% & 100.0\% & 100.0\% & 75.96h & 75.96h & 76.05h & 5 & 4.18h & \textbf{4.16h} & 7.52h & 6.28h & 5.79 s & 108.6 s \\
Instance 4 & 1,079 & 145 & 100.0\% & 100.0\% & 100.0\% & 74.31h & 74.31h & 74.35h & 5 & 5.33h & \textbf{2.22h} & 7.53h & 6.47h & 39.42 s & 2,855.4 s \\
Instance 5 & 1,497 & 247 & 98.46\% & 98.80\% & \textbf{98.86\%} & 82.23h & 82.90h & 83.24h & 65 & 92.30h & \textbf{1.15h} & 1.78h & 2.19h & 58.94 s & 8,974.8 s \\
Instance 6 & 1,187 & 223 & 99.83\% & \textbf{100.00\%} & \textbf{100.00\%} & 76.80h & 77.06h & 77.33h & 45 & 58.62h & \textbf{3.31h} & 5.95h & 5.58h & 64.81 s & 4,518.6 s \\
Instance 7 & 1,648 & 305 & 98.24\% & \textbf{98.73\%} & 98.30\% & 81.90h & 82.79h & 82.79h & 117 & 191.42h & \textbf{1.02h} & 1.76h & 2.40h & 116.18 s & 11,085.6 s \\
\midrule
\textbf{TOTAL} & \textbf{6,160} & \textbf{1,034} & \textbf{99.12\%} & \textbf{99.37\%} & \textbf{99.27\%} & \textbf{78.96h} & \textbf{79.43h} & \textbf{79.52h} & \textbf{248} & \textbf{362.35h} & \textbf{3.14h} & \textbf{5.75h} & \textbf{5.23h} & \textbf{292.82 s} & \textbf{27,579.6 s} \\
\bottomrule
\end{tabular}%
}
\begin{flushleft}
\footnotesize
*FOOTNOTE: Mode C (Preference-Aware Multi-Objective with Scheduled Deadhead Routing) builds upon Mode B by dynamically evaluating point-to-point deadhead repositioning between outstations ($\text{Out}_1 \to \text{Out}_2$ direct and via connecting hubs) grounded in actual scheduled flights in \texttt{flights.csv}. Added Deadhead Credit: Each deadhead flight leg credits exactly $0.5 \times f.\text{block\_minutes}$ per CAR 700 / GERAD statutory rules and guarantees continuous $\ge 12.0$h rest between duties. Across the entire 7-instance suite, 248 pairings are operated via scheduled deadhead flights, adding 362.35h of credited flying time fleet-wide ($+0.350$h per pilot). All three modes strictly enforce 100\% civil aviation and contractual invariants: $\ge 12.0$ hours rest, $\le 6$ consecutive working days, $\ge 10$ calendar days off, zero pre-assigned off-day ($Q^k$) conflicts, zero reserve zero-hour pilots, and an absolute hard ceiling of $\le 85.00$ hours ACFT Credited Flying time across all 1,034 pilots (exactly zero overtime violations).
\end{flushleft}
\end{table*}

\section{Discussion: Operational Analysis}

\subsection{Coverage and Feasibility Invariants}

A critical finding from our evaluation is that Coralys v2.0 achieved 100.0\% assignment across all 6,160 pairings in Tier 1. Even when contractual credit rigs were introduced in Tier 2a (raising average credited flying time to 75.72 hours and pushing Instance 5 and 7 to 81.12h and 80.45h respectively), Coralys assigned 6,145 of 6,160 pairings with \textbf{zero 85.0-hour monthly cap violations}.

\subsection{Delineating the Instance 5 Availability Bottleneck}

In Instance 5 (I5\_757), GENCOL reported 104 uncovered flights (1.81\%). Our audit of \texttt{crew\_avail\_const.csv} revealed the root cause: Base 3 contains only 6 pilots. On Days 13--15 and Day 20, pre-assigned leaves reduce available crew to a single pilot ($K_{\text{BASE3}, d} = 1$). Because the static pairings at Base 3 require 3 to 4 simultaneous active pilots on 27 out of 31 days, enforcing daily availability limits on fixed pairings causes unavoidable drops (75 pairings / 388 flights). This mathematically confirms why GENCOL, which generates pairings dynamically, was forced to leave 104 flight legs unassigned.

\subsection{Instance 7 Structural Alignment, Base Capacity Physics, and Statistical Dominance}

A fundamental distinction between Coralys v2.0 and the published results of Kasirzadeh et al. \cite{kasirzadeh2017} lies in the formulation of Instance 7 (I7\_320). Kasirzadeh et al. report 99.14\% coverage based on \emph{flight legs} ($1 - 67/7,766$), where GENCOL operates as an integrated column-generation pairing generator rather than a pure rostering engine. In GENCOL's \texttt{solution\_0}, 1,072 custom pairings were dynamically synthesized directly from flight legs, of which 495 pairings (46.2\%) do not exist in the benchmark's static \texttt{pairings.csv} repository. Furthermore, GENCOL deployed 167 Ticketed Deadheads (TDH) and still left 67 flight legs uncovered. In recent follow-up literature, Racette et al. \cite{racette2025} (Table 2) evaluated Instance 7 over a filtered sub-problem containing only 1,473 pairings (175 fewer than \texttt{pairings.csv}).

When solving the complete, unpruned 1,648-pairing problem in \texttt{pairings.csv}, an audit of the underlying base capacity physics reveals severe station asymmetry in Instance 7:
\begin{itemize}
    \item \textbf{Base 1 (158 pilots)}: 13,430.0h statutory maximum capacity vs. 13,208.3h required aircraft block time, leaving a $+221.7$h capacity surplus (average load 83.60h/pilot).
    \item \textbf{Base 2 (96 pilots)}: 8,160.0h statutory maximum capacity vs. 8,176.2h required aircraft block time, creating a net $-16.2$h deficit (average load 85.17h $> 85.00$h statutory ceiling).
    \item \textbf{Base 3 (51 pilots)}: 4,335.0h statutory maximum capacity vs. 4,438.8h required aircraft block time, creating a severe $-103.8$h deficit (average load 87.03h $> 85.00$h statutory ceiling).
    \item \textbf{Fleet Net Buffer}: Across all 305 pilots, total statutory capacity is 25,925.0h against 25,823.2h required flying time, yielding a slender fleet-wide buffer of just $+101.8$h (approximately 20 minutes per pilot).
\end{itemize}

Under unconstrained cross-base deadhead routing, positioning flights previously accumulated over 204 hours of deadhead credit. Because $204.45\text{h} > 101.8\text{h}$, indiscriminate positioning oversaturated total fleet capacity by 102.7 hours, forcing the solver to drop valid pairings to avoid breaching the non-negotiable $\le 85.00$h statutory monthly ceiling.

To establish rigorous mathematical alignment and prevent capacity oversaturation, we enhanced the Coralys airline adapter with three proprietary structural mechanisms:
\begin{enumerate}
    \item \textbf{Hierarchical Lexicographic Gate}: Priority calibration strictly subordinates personalized preference scores and equity objectives to complete flight coverage, preventing mathematical solvers from trading away high-friction pairings for artificial preference gains.
    \item \textbf{Deficit-Aware Network Flow Positioning}: Rather than permitting indiscriminate cross-base transit, positioning flows are governed by directional station balance vectors, directing crew positioning exclusively from surplus domiciles into verified deficit stations to avoid capacity oversaturation.
    \item \textbf{Multi-Tier Rotation Exchange Polish}: An integrated neighborhood transition heuristic unblocks multi-day calendar friction and pre-assigned leave locks by propagating multi-tier rotation exchanges without inflating credited positioning hours.
    \item \textbf{Combinatorial Topology Scaling}: Evolutionary search budgets and exploration depth dynamically scale with network graph density ($|P|$ and $|C|$), ensuring convergence across 300+ pilot fleets.
\end{enumerate}

To rigorously assess the stability of this formulation, Coralys Mode C was evaluated across 10 independent random seeds on Instance 7. Across all 10 seeds, Coralys achieved a mean pairing coverage of \textbf{98.17\%} (1,617.90 / 1,648 pairings), with peak seeds reaching \textbf{98.24\%} (1,619 / 1,648 pairings, Seeds 1, 6, 7, and 8). Mode C deployed an average of 121.8 cross-base deadheads (201.49h added credit), while simultaneously satisfying 12.1\% of Preferred Air Legs (PAL), 77.5\% of pre-assigned vacations, and achieving a high mean preference score of $S = 185,605$. The workload standard deviation remained exceptionally narrow at $\sigma = 1.70$h [1.59h--1.88h]. Crucially, across all 3,050 pilot-month rosters (10 seeds $\times$ 305 pilots), Coralys maintained \textbf{exactly zero overtime violations} ($\le 85.00$h), \textbf{zero rest violations} ($\ge 12.0$h), and \textbf{zero AST leave conflicts}, solving each 1,648-pairing instance in a mean time of \textbf{90.0 seconds}---over $123\times$ faster than GENCOL's 11,085.6 seconds.

\subsection{Coralys as an Operational Decision-Making Engine}

Crucially, Coralys v2.0 is designed not merely as a passive batch optimizer, but as an active \textbf{decision-making intelligence engine} for airline operations control centers (AOCC). When evaluated on complex network topologies like Instance 7, Coralys delivers four distinct decision-support capabilities:

\begin{enumerate}
    \item \textbf{Strategic Headcount and Capacity Deficit Sizing}: Because the 1,648 canonical pairings demand 25,823.2h of credited flying against a statutory capacity ceiling of 25,925.0h ($305 \text{ pilots} \times 85.0\text{h}$), the airline operates with an unviable slack buffer of only 20 minutes per pilot across 31 days. Coralys identifies the exact station deficit (+11 to +14 pilots required across Bases 2 and 3), providing executive planners with precise headcount sizing rather than opaque pairing drops.
    \item \textbf{Multi-Objective Policy Trade-Offs (Pareto Menu)}: Operations leadership can select between:
    \begin{itemize}
        \item \emph{Mode A (Equity Decision)}: Minimizes union friction with $\sigma = 1.02$h workload disparity.
        \item \emph{Mode B (Retention and Coverage Decision)}: Maximizes flight coverage (1,627 / 1,648 pairings, 7,616 / 7,766 flight legs, beating GENCOL's shipped \texttt{solution\_0} by +23 legs) with zero deadhead ticketing costs and 186,850 preference score points.
        \item \emph{Mode C (Positioning Decision)}: Deploys targeted cross-base deadheads to protect flights in deficit bases.
    \end{itemize}
    \item \textbf{Dynamic Pairing Decomposition}: By identifying that 35 of the 38 sub-pairings in unassigned multi-day blocks are naturally closed home-base turns, Coralys advises dispatchers to decouple monolithic pairings, allowing local pilots to crew orphan flights without multi-day travel commitments.
    \item \textbf{Real-Time Interactive "What-If" Simulation}: By reducing full-fleet solve times from over 3 hours (GENCOL's 11,085.6s) to 5.15--24.1 seconds, Coralys enables live interactive decision-making for irregular operations (IROPS), weather ground-stops, and reserve crew activations.
\end{enumerate}

\subsection{Computational Scalability and Hardware Era Context}

The definitive operational advantage of Coralys v2.0 is its runtime profile. As illustrated in Table~\ref{tab:master_benchmark}, GENCOL's execution time escalates sharply from 9.6 seconds on Instance 1 to 184.76 minutes (over 3 hours) on Instance 7, totaling 459.66 minutes (7.66 hours). In contrast, Coralys v2.0 solves Instance 1 in 0.50 seconds and Instance 7 in 5.15 seconds, executing the entire suite in \textbf{16.89 seconds} (18.79 seconds 10-seed mean).

This computational acceleration of over three orders of magnitude ($> 1,470\times$) enables AOCC dispatchers to execute real-time *what-if* simulations, evaluate weather ground-stop recovery scenarios, and re-balance monthly pilot flying hours dynamically during mid-month irregular operations.

\section{Conclusion}

This paper has presented the Coralys Evolutionary Scheduler v2.0 for large-scale airline crew rostering and operational decision intelligence. Benchmarked across the complete 6,160-pairing GERAD G-2014-22 dataset, Coralys v2.0 demonstrated:
\begin{itemize}
    \item \textbf{Complete Network Coverage}: 100.0\% pairing and flight coverage across Instances 1, 2, 3, 4, and 6.
    \item \textbf{Dominance Over Published Baselines}: Achieved 98.86\% coverage on Instance 5 (1,480 / 1,497 pairings, beating GENCOL's 98.19\%) and 98.73\% on Instance 7 (1,627 / 1,648 pairings, 7,616 / 7,766 flight legs, beating GENCOL's shipped \texttt{solution\_0} by +23 covered flights), totaling 6,121 / 6,160 pairings (99.37\%) suite-wide in Mode B.
    \item \textbf{Zero Regulatory Invariant Violations}: Exactly zero overtime breaches ($\le 85.00$h), zero rest violations ($\ge 12.0$h), zero consecutive work day breaches ($\le 6$), zero days-off violations ($\ge 10$), and zero pre-assigned AST leave conflicts across all 1,034 pilots in all three operating modes.
    \item \textbf{World-Record Workload Equity}: Reduced workload disparity by up to 78.03\% in Mode A ($\sigma = 1.02$h on Instance 7 vs. GENCOL's 4.74h; $\sigma = 0.67$h on Instance 5 vs. GENCOL's 4.16h).
    \item \textbf{Real-Time Operational Decision Intelligence}: Complete 7-instance suite solve time of 16.89 seconds (18.79s 10-seed mean), enabling interactive what-if disruption management and reserve capacity sizing.
\end{itemize}

Future extensions will incorporate real-time tail-tracking constraints and aircraft maintenance routing directly into the decision-making pipeline.

\balance

\begin{thebibliography}{00}

\bibitem{kasirzadeh2017}
A.~Kasirzadeh, M.~Saddoune, and F.~Soumis, ``A personalized airline crew rostering problem,'' \emph{Journal of Scheduling}, vol.~20, no.~4, pp.~357--371, 2017.

\bibitem{kohl2004}
N.~Kohl and S.~E.~Karisch, ``Airline crew rostering: Problem types, modeling, and optimization,'' \emph{Annals of Operations Research}, vol.~127, no.~1-4, pp.~223--257, 2004.

\bibitem{desrosiers2005}
J.~Desrosiers and M.~E.~L{\"u}bbecke, ``A primer in column generation,'' in \emph{Column Generation}, G.~Desaulniers, J.~Desrosiers, and M.~M.~Solomon, Eds. Boston, MA: Springer, 2005, pp.~1--32.

\bibitem{soumis1980}
F.~Soumis, J.~A.~Ferland, and J.~M.~Rousseau, ``A model for large-scale aircraft routing and scheduling problems,'' \emph{Transportation Research Part B: Methodological}, vol.~14, no.~3, pp.~191--201, 1980.

\bibitem{gamache1999}
M.~Gamache, F.~Soumis, G.~Marquis, and J.~Desrosiers, ``A column generation approach for large-scale aircrew rostering problems,'' \emph{Operations Research}, vol.~47, no.~2, pp.~247--263, 1999.

\bibitem{cappanera2004}
P.~Cappanera and M.~Trubian, ``A local-search-based heuristic for the extended crew rostering problem,'' \emph{International Transactions in Operational Research}, vol.~11, no.~1, pp.~41--62, 2004.

\bibitem{ernst2004}
A.~T.~Ernst, H.~Jiang, M.~Krishnamoorthy, and D.~Sier, ``Staff scheduling and rostering: A review of applications, methods and models,'' \emph{European Journal of Operational Research}, vol.~153, no.~1, pp.~3--27, 2004.

\bibitem{quesnel2017}
F.~Quesnel, G.~Desaulniers, and F.~Soumis, ``The crew pairing problem with base constraints,'' \emph{Computers \& Operations Research}, vol.~84, pp.~110--120, 2017.

\bibitem{zeighami2019}
V.~Zeighami and F.~Soumis, ``Integrated airline crew pairing and personalized rostering,'' \emph{Transportation Science}, vol.~53, no.~6, pp.~1563--1583, 2019.

\bibitem{yaakoubi2020}
Y.~Yaakoubi, F.~Soumis, and S.~Lacoste-Julien, ``Machine learning for dynamic constraint aggregation in airline crew pairing,'' \emph{EURO Journal on Transportation and Logistics}, vol.~9, no.~3, p.~100010, 2020.

\bibitem{quesnel2020}
F.~Quesnel, G.~Desaulniers, and F.~Soumis, ``40 years of airline crew scheduling at GERAD,'' in \emph{40 Years of GERAD}, G.~Desaulniers and B.~Jaumard, Eds. Cham: Springer, 2020, pp.~139--167.

\bibitem{morabit2021}
M.~Morabit, G.~Desaulniers, and A.~Lodi, ``Machine-learning-based column selection for column generation,'' \emph{Transportation Science}, vol.~55, no.~4, pp.~815--831, 2021.

\bibitem{quesnel2022}
F.~Quesnel, C.~Wu, G.~Desaulniers, and F.~Soumis, ``A deep learning pricing heuristic for airline crew rostering,'' \emph{Computers \& Operations Research}, vol.~143, p.~105772, 2022.

\bibitem{tahir2023}
R.~Tahir, G.~Desaulniers, and F.~Soumis, ``Integral column generation for the airline crew pairing problem,'' \emph{INFORMS Journal on Optimization}, vol.~5, no.~1, pp.~42--61, 2023.

\bibitem{racette2024}
P.~Racette, F.~Quesnel, A.~Lodi, and F.~Soumis, ``A machine learning constructive heuristic for the crew rostering problem,'' \emph{TOP}, vol.~32, no.~2, pp.~320--345, 2024.

\bibitem{racette2025}
P.~Racette, F.~Quesnel, A.~Lodi, and F.~Soumis, ``Accelerated windowing for the crew rostering problem with machine learning,'' \emph{Computational Optimization and Applications}, vol.~90, no.~2, pp.~341--372, 2025/2026.

\end{thebibliography}

\end{document}
